% gerado por ferramentas/md2tex_enunciado.py a partir de questoes/q111.md (não editar à mão)
A busca por fontes alternativas de energia é uma das grandes preocupações da humanidade devido ao possível esgotamento dos combustíveis fósseis e à degradação ambiental causada por eles. Como a demanda de energia cresce significativamente, há um grande investimento em outros tipos de geração de energia, tais como hidrelétrica, nuclear, solar, termelétrica a biomassa e eólica.\par
\fonte{BIZERRA, Ayla M. C.; QUEIROZ, Jorge L. A.; COUTINHO, Demétrios A. M. O impacto ambiental dos combustíveis fósseis e dos biocombustíveis: as concepções de estudantes do ensino médio sobre o tema. \textbf{Revbea}, n. 3, 2018 (adaptado).}
\figq[0.78\linewidth]{q111_fig1.png}
Dentre as fontes de energia citadas, o tipo de energia oriunda de uma fonte não renovável, que independe das condições climáticas para ser obtida, é a\par
\begin{alternativas}
  \item nuclear, gerada pela fissão nuclear do urânio.
  \item hidrelétrica, gerada pela movimentação da água.
  \item solar, gerada pelo uso de células fotovoltaicas.
  \item termelétrica, produzida pela combustão de biomassa.
  \item eólica, produzida pelo movimento das pás de aerogeradores.
\end{alternativas}
